% sections/discussion.tex --- Discussion
\section{Discussion}
\label{sec:discussion}

We use this section to draw out the methodological consequences of
the results in Section~\ref{sec:evaluation}, to call out the most
pressing threats to validity, and to position \bench{} in the
broader research economy of audit analytics.

\subsection{What baselines learn, and what they do not}
\label{sec:discussion:learn}

The expected gap between the temporal and typology splits
(Section~\ref{sec:evaluation:cross}) is the single most
methodologically suggestive feature of the benchmark.  Temporal
generalisation is comparatively easy: held-out months draw from the
same anomaly categories that the model has already seen during
training, and the relevant signal is largely \emph{within
distribution}.  Typology generalisation is harder by design: the
held-out anomaly categories are absent from the training pool, and
the model must infer that a posting is anomalous on the basis of
\emph{what is normal}, not on the basis of what previously labelled
positives looked like.

Strong GNN baselines are, in our preliminary experience, much more
sensitive to this distinction than tabular baselines.  GNNs that
do well on the temporal split frequently do not preserve their
margin on the typology split, suggesting that they are exploiting
class-specific structural cues rather than building a robust
representation of normality.  This is a candidate driver for
methodological work on \emph{anomaly definitions} that do not rely
on fully supervised typology coverage---auto-encoders trained on
clean training pools, contrastive losses that pair postings against
their document-chain predecessors, or open-set detectors that
explicitly admit unseen categories at test time.

\subsection{Threats to validity}
\label{sec:discussion:threats}

\textbf{Synthetic-to-real transfer.}  The most serious threat to
the validity of any synthetic benchmark is that methods which
dominate on synthetic data fail to transfer to real client
ledgers.  \bench{} is calibrated against a $364$M-posting empirical
study (Section~\ref{sec:empirical}), and the released configuration
reproduces the line-item, account-population, and temporal
distributions reported there.  This calibration is necessary but
not sufficient: distributional fidelity at the marginal level does
not guarantee fidelity at the joint level, and we make no claim
that a method's \bench{} score directly predicts its score on a
client ledger.  The role of \bench{} is to provide a
\emph{reference frame} on which methods can be compared,
ablated, and reproduced; the validation against client data
remains a separate exercise that practitioners must conduct
themselves and under their own confidentiality protocols.

\textbf{Anomaly definitions.}  The labels in
\texttt{labels/anomaly\_labels.csv} are determined by the
generator's anomaly-injection module
(Section~\ref{sec:datasynth:anomaly}); they reflect the categories
the module is configured to inject and the propagation rule from
entries to edges.  A method that targets a different operational
definition of ``anomaly''---for example, a definition centred on
control-failure rather than on fraud---should expect to see lower
numbers on \bench{} than its target definition would suggest, and
should report this gap explicitly when claiming applicability.

\textbf{Single industry profile.}  The released v$1.0$ configuration
is manufacturing-only.  Cross-industry generalisation is not
something the v$1.0$ benchmark numbers can measure.  The
configuration interface supports the four other industry profiles
shipped with \system{}, and we expect to release multi-industry
splits in v$1.1$.

\subsection{Confidentiality, reproducibility, and the role of public
benchmarks in audit research}
\label{sec:discussion:reproducibility}

The principal motivation for \bench{} is the observation
(Section~\ref{sec:introduction:gap}) that a research community
cannot accumulate methodological progress on private,
irreproducible datasets.  This observation cuts in two directions
that are worth separating.

First, public benchmarks do not displace empirical work on real
client data.  An audit method that wins a leaderboard on \bench{}
is not, by that fact alone, ready for production deployment; it
must still be validated against the client ledger it will run
against, under the auditor's own quality controls.  A reference
benchmark is a coordination device that lets methods be compared
and ablated; it is not a substitute for client engagement.

Second, public benchmarks lower the cost of method development for
audit practitioners.  An audit team that wishes to evaluate a new
fraud-detection technique against the published state of the art
no longer needs to license and replicate a corpus or to negotiate
data-sharing agreements; the team can run the technique against
\bench{} and read the leaderboard.  This is a particularly large
gain for smaller audit firms, in-house audit functions, and
academic researchers, whose access to client data is structurally
more restricted than that of the largest firms.

\subsection{Implications for practitioners}
\label{sec:discussion:implications}

For practitioners specifically: \bench{} is a measuring stick.
It does not---and is not designed to---tell a CAE or audit
methodology lead which detection technique to deploy on next
year's engagement.  It does provide a defensible answer to the
question ``how does technique X compare to the state of the art
on equivalent data?''  When a vendor claims that a model achieves
$\AUCPR = 0.X$, the practitioner can now ask whether the same
model achieves the same $\AUCPR$ on \bench{}, and whether the
model's performance on the typology-hold-out split is comparable.
A negative answer is informative.
